% ============================================================
%  USPSC-Cap4-Avaliacao_Experimental.tex
%  Capítulo 4 — Avaliação Experimental
%  TCC: Análise de Sentimento de Publicações do X/Twitter
%       no Mercado Financeiro Brasileiro com FinBERT-PT-BR
% ============================================================
\chapter{Avaliação Experimental} \label{Cap:Ava_Exp}

Este capítulo apresenta os resultados da avaliação comparativa das quatro abordagens de classificação de sentimento descritas na \autoref{sec:abordagens}: OpLexicon~v3.0, SentiLex-PT, FinBERT-PT-BR e BERTimbau ajustado. A avaliação foi conduzida exclusivamente sobre o conjunto de teste \textit{hold-out} definido na \autoref{subsec:particionamento}, composto por 373 publicações rotuladas manualmente pelo pesquisador, em conformidade com o protocolo detalhado no \autoref{Cap:Metodologia}.

A \autoref{sec:config_experimental} descreve as condições de configuração do experimento. A \autoref{sec:resultados_abordagem} apresenta os resultados individuais de cada abordagem. A \autoref{sec:analise_comparativa} realiza a análise comparativa consolidada entre as abordagens e, por fim, a \autoref{sec:discussao_resultados} interpreta os achados à luz da literatura e dos objetivos da pesquisa.

% ============================================================
%  USPSC-Cap4-Avaliacao_Experimental_4.1-Configuracao_Experimental.tex
%  Seção 4.1 — Configuração Experimental
% ============================================================

\section{Configuração Experimental}
\label{sec:config_experimental}

\subsection{Conjunto de Teste}
\label{subsec:conjunto_teste}

A avaliação comparativa foi realizada sobre um conjunto de teste composto por 373 publicações, extraídas do subconjunto rotulado manualmente conforme o protocolo de anotação descrito na \autoref{sec:rotulagem}. O conjunto corresponde à partição de teste (15\%) obtida por particionamento estratificado com \texttt{random\_state=42}, conforme detalhado na \autoref{subsec:particionamento}. A estratificação preservou a distribuição original das três classes de sentimento no corpus rotulado, resultando na composição apresentada na \autoref{tab:distribuicao_teste}.

\begin{table}[!htbp]
	\caption[Distribuição de classes no conjunto de teste]{Distribuição de classes no conjunto de teste ($n = 373$)}
	\label{tab:distribuicao_teste}
	\centering
	\begin{tabular}{lcc}
		\hline
		\textbf{Classe} & \textbf{Instâncias} & \textbf{Proporção} \\
		\hline
		Positivo  & 59  & 15,8\% \\
		Neutro    & 266 & 71,3\% \\
		Negativo  & 48  & 12,9\% \\
		\hline
		\textbf{Total} & \textbf{373} & \textbf{100,0\%} \\
		\hline
	\end{tabular}
	\legend{Fonte: Elaborada pelo autor.}
\end{table}

A predominância da classe neutra (71,3\%) é consistente com o caráter editorial e factual das publicações dos veículos analisados, conforme discutido na \autoref{subsec:lim_representatividade}, e constitui um fator relevante na interpretação das métricas de cada abordagem, em particular do F1-score macro, que é calculado sem ponderação pelo tamanho de cada classe.

\subsection{Protocolo de Avaliação}
\label{subsec:protocolo_avaliacao}

Todas as quatro abordagens avaliadas — OpLexicon~v3.0, SentiLex-PT, FinBERT-PT-BR e BERTimbau ajustado --- foram aplicadas exatamente ao mesmo conjunto de 373 instâncias de teste, sem que qualquer abordagem tenha tido acesso a esses exemplos durante o desenvolvimento ou o treinamento. Para o BERTimbau ajustado, o particionamento estratificado garantiu que os exemplos de teste nunca fizeram parte das partições de treinamento (70\%) ou de validação (15\%) utilizadas no \textit{fine-tuning}.

As abordagens léxicas (OpLexicon e SentiLex-PT) e o FinBERT-PT-BR operam de forma determinística e não consomem dados de treinamento do corpus local; portanto, sua avaliação no conjunto de teste é direta e não sujeita a viés de vazamento (\textit{data leakage}). As métricas calculadas para cada abordagem são as definidas na \autoref{subsec:metricas_avaliacao}: acurácia, precisão, revocação e F1-score por classe, F1-score macro e matriz de confusão. A implementação do protocolo de avaliação está centralizada no módulo \texttt{app/core/evaluation/metrics.py}, invocado via a função \texttt{evaluate(classificator)}.

\subsection{Configuração do Fine-Tuning do BERTimbau}
\label{subsec:config_finetuning}

O processo de \textit{fine-tuning} do BERTimbau, descrito em detalhes na \autoref{subsubsec:finetuning_bertimbau}, foi executado com os hiperparâmetros consolidados no \autoref{qua:hiperparametros_bertimbau}.


\begin{quadro}[!htbp]
	\caption[Hiperparâmetros do fine-tuning do BERTimbau]{Hiperparâmetros utilizados no \textit{fine-tuning} do BERTimbau}
	\label{qua:hiperparametros_bertimbau}
	\centering
	\begin{tabular}{ll}
		\hline
		\textbf{Hiperparâmetro} & \textbf{Valor} \\
		\hline
		Modelo base                          & \texttt{neuralmind/bert-base-portuguese-cased} \\
		Taxa de aprendizado                  & $5 \times 10^{-5}$ \\
		Tamanho do batch                     & 8 \\
		Número máximo de épocas              & 12 \\
		Comprimento máximo (\textit{tokens}) & 512 \\
		\textit{Warmup ratio}                & 0,10 \\
		\textit{Weight decay}                & 0,01 \\
		\textit{Early stopping}              & Paciência de 2 épocas (F1 macro na dobra de validação) \\
		Validação cruzada                    & 4 dobras estratificadas (\textit{best fold}: 1, F1 macro val. = 90,00\%) \\
		Função de perda                      & \texttt{CrossEntropyLoss} com pesos de classe balanceados \\
		Semente aleatória                    & 42 \\
		\hline
	\end{tabular}
	\legend{Fonte: Elaborada pelo autor.}
\end{quadro}


O desbalanceamento de classes foi mitigado por meio de uma subclasse personalizada do \texttt{Trainer} da biblioteca \textit{HuggingFace Transformers} \cite{wolf2020transformers} (\texttt{WeightedTrainer}), que recalcula a função de perda com pesos inversamente proporcionais à frequência de cada classe, calculados via \texttt{compute\_class\_weight('balanced')} do scikit-learn. O treinamento foi executado sobre a partição de 70\% do corpus rotulado, com monitoramento sobre a partição de validação (15\%) para ativação do \textit{early stopping}.

\subsection{Ambiente Computacional}
\label{subsec:ambiente}

O \textit{pipeline} de avaliação foi executado em ambiente Python 3.11, com as principais dependências listadas no \autoref{qua:dependencias}.

\begin{quadro}[!htbp]
	\caption[Principais dependências do ambiente de execução]{Principais dependências do ambiente de execução}
	\label{qua:dependencias}
	\centering
	\begin{tabular}{ll}
		\hline
		\textbf{Biblioteca} & \textbf{Finalidade} \\
		\hline
		\texttt{transformers} (HuggingFace) & Carregamento e inferência dos modelos BERT \\
		\texttt{scikit-learn}               & Métricas de avaliação e pesos de classe \\
		\texttt{torch} (PyTorch)            & Backend de treinamento e inferência neural \\
		\texttt{pandas} / \texttt{numpy}    & Manipulação de dados e matrizes \\
		\texttt{psycopg2}                   & Acesso ao banco de dados PostgreSQL \\
		\texttt{loguru}                     & Registro de eventos do \textit{pipeline} \\
		\hline
	\end{tabular}
	\legend{Fonte: Elaborada pelo autor.}
\end{quadro}

O banco de dados PostgreSQL foi provisionado em contêiner Docker, e os artefatos do \textit{fine-tuning} do BERTimbau --- incluindo a matriz de confusão e o relatório de classificação do conjunto de teste --- foram persistidos no diretório \texttt{models/bert-timbau-sentiment/eval/} do repositório do \textit{pipeline}.
% ============================================================
%  USPSC-Cap4-Avaliacao_Experimental_4.2-Resultados_por_Abordagem.tex
%  Seção 4.2 — Resultados por Abordagem
% ============================================================

\section{Resultados por Abordagem}
\label{sec:resultados_abordagem}

Esta seção apresenta os resultados individuais de cada abordagem de classificação sobre o conjunto de teste ($n = 373$). Para cada abordagem são reportados: o relatório completo de classificação por classe (\autoref{subsec:metricas_avaliacao}) e a matriz de confusão, que evidencia os padrões sistemáticos de acerto e erro.


% ── 4.2.1 OpLexicon ──────────────────────────────────────────────────────────

\subsection{OpLexicon v3.0}
\label{subsec:res_oplexicon}

O OpLexicon~v3.0 \cite{souza2011oplexicon} classifica cada publicação pela polaridade líquida dos termos presentes no léxico, após a etapa de pré-processamento. O resultado apresentado na \autoref{tab:metrics_oplexicon} revela o desempenho mais baixo entre as quatro abordagens avaliadas.


\begin{table}[!htbp]
	\caption[Métricas de desempenho — OpLexicon v3.0]{Métricas de desempenho do OpLexicon v3.0 ($n = 373$)}
	\label{tab:metrics_oplexicon}
	\centering
	\begin{tabular}{lcccc}
		\hline
		\textbf{Classe}    & \textbf{Precisão} & \textbf{Revocação} & \textbf{F1-Score} & \textbf{Suporte} \\
		\hline		
		Positivo           & 0,18              & 0,68               & 0,29              & 59  \\
		\hline
		Neutro             & 0,81              & 0,13               & 0,22              & 266 \\
		\hline
		Negativo           & 0,20              & 0,46               & 0,28              & 48  \\
		\hline
		\textit{Macro avg} & 0,40              & 0,42               & 0,26              & 373 \\
		\hline
		\textit{Acurácia}  & ---               & ---                & 0,26              & 373 \\
		\hline
	\end{tabular}
	\legend{Fonte: Elaborada pelo autor.}
\end{table}


\begin{table}[!htbp]
	\caption[Matriz de confusão — OpLexicon v3.0]{Matriz de confusão do OpLexicon v3.0. Linhas: rótulo verdadeiro; colunas: rótulo predito.}
	\label{tab:cm_oplexicon}
	\centering
	\begin{tabular}{lccc}
		\hline
		& \textbf{Pred. Positivo} & \textbf{Pred. Neutro} & \textbf{Pred. Negativo} \\
		\hline
		\textbf{Real Positivo} & 40                      & 4                     & 15                      \\
		\hline
		\textbf{Real Neutro}   & 157                     & 34                    & 75                      \\
		\hline
		\textbf{Real Negativo} & 22                      & 4                     & 22                      \\
		\hline
	\end{tabular}
	\legend{Fonte: Elaborada pelo autor.}
\end{table}

O OpLexicon alcançou acurácia de 25,74\% e F1-score macro de 26,12\%, indicando desempenho sistematicamente inferior ao classificador aleatório estratificado. A matriz de confusão (\autoref{tab:cm_oplexicon}) revela o principal padrão de falha: 157 das 266 publicações neutras (59,0\%) foram classificadas como positivo, evidenciando uma tendência de atribuir polaridade positiva a publicações jornalísticas factuais cujo vocabulário financeiro não carrega sinal negativo explícito. A classe neutro, dominante no corpus, obteve revocação de apenas 13\%, o que significa que a abordagem léxica é incapaz de identificar o padrão de ausência de sentimento predominante neste tipo de dado.

% ── 4.2.2 SentiLex-PT ────────────────────────────────────────────────────────

\subsection{SentiLex-PT}
\label{subsec:res_sentilex}

O SentiLex-PT \cite{carvalho2015sentilex} é um léxico de sentimentos para o português construído a partir de entradas morfologicamente anotadas. Os resultados são apresentados na \autoref{tab:metrics_sentilex}.

\begin{table}[!htbp]
	\caption[Métricas de desempenho — SentiLex-PT]{Métricas de desempenho do SentiLex-PT ($n = 373$)}
	\label{tab:metrics_sentilex}
	\centering
	\begin{tabular}{lcccc}
		\hline
		\textbf{Classe}    & \textbf{Precisão} & \textbf{Revocação} & \textbf{F1-Score} & \textbf{Suporte} \\
		\hline
		Positivo           & 0,19              & 0,59               & 0,29              & 59  \\
		\hline
		Neutro             & 0,78              & 0,20               & 0,31              & 266 \\
		\hline
		Negativo           & 0,18              & 0,48               & 0,27              & 48  \\
		\hline
		\textit{Macro avg} & 0,38              & 0,42               & 0,29              & 373 \\
		\hline
		\textit{Acurácia}  & ---               & ---                & 0,29              & 373 \\
		\hline
	\end{tabular}
	\legend{Fonte: Elaborada pelo autor.}
\end{table}

\begin{table}[!htbp]
	\caption[Matriz de confusão — SentiLex-PT]{Matriz de confusão do SentiLex-PT. Linhas: rótulo verdadeiro; colunas: rótulo predito.}
	\label{tab:cm_sentilex}
	\centering
	\begin{tabular}{lccc}
		\hline
		& \textbf{Pred. Positivo} & \textbf{Pred. Neutro} & \textbf{Pred. Negativo} \\
		\hline
		\textbf{Real Positivo} & 35                      & 7                     & 17                      \\
		\hline
		\textbf{Real Neutro}   & 129                     & 52                    & 85                      \\
		\hline
		\textbf{Real Negativo} & 17                      & 8                     & 23                      \\
		\hline
	\end{tabular}
	\legend{Fonte: Elaborada pelo autor.}
\end{table}

O SentiLex-PT obteve acurácia de 29,49\% e F1-score macro de 29,00\%, ligeiramente superior ao OpLexicon, mas ainda muito abaixo de um classificador por maioria de classe (\textit{majority class baseline}). A matriz de confusão (\autoref{tab:cm_sentilex}) evidencia que 129 das 266 publicações neutras (48,5\%) foram classificadas como positivo, padrão análogo ao observado no OpLexicon. A revocação da classe neutra é igualmente baixa (20\%), confirmando a dificuldade estrutural de abordagens léxicas em reconhecer ausência de polaridade em textos de caráter informativo. O SentiLex-PT distribui mais erros entre positivo e negativo do que o OpLexicon, refletindo diferenças na cobertura léxica de cada recurso.

% ── 4.2.3 FinBERT-PT-BR ──────────────────────────────────────────────────────

\subsection{FinBERT-PT-BR}
\label{subsec:res_finbert}

O FinBERT-PT-BR \cite{santos2023finbert} é um modelo BERT especializado para o domínio financeiro em língua portuguesa, pré-treinado sobre corpus de notícias financeiras. Neste trabalho, o modelo é utilizado em modo de inferência direta (\textit{zero-shot} em relação ao corpus local, sem \textit{fine-tuning} sobre as publicações coletadas), precedido apenas da etapa de normalização de espaçamento descrita na \autoref{subsec:fluxo_finbert}, conforme detalhado na \autoref{subsubsec:protocolo_inferencia}. Os resultados são apresentados na \autoref{tab:metrics_finbert}.

\begin{table}[!htbp]
	\caption[Métricas de desempenho — FinBERT-PT-BR]{Métricas de desempenho do FinBERT-PT-BR ($n = 373$)}
	\label{tab:metrics_finbert}
	\centering
	\begin{tabular}{lcccc}
		\hline
		\textbf{Classe}    & \textbf{Precisão} & \textbf{Revocação} & \textbf{F1-Score} & \textbf{Suporte} \\
		\hline
		Positivo           & 0,41              & 0,47               & 0,44              & 59  \\
		\hline
		Neutro             & 0,88              & 0,30               & 0,45              & 266 \\
		\hline
		Negativo           & 0,20              & 0,88               & 0,32              & 48  \\
		\hline
		\textit{Macro avg} & 0,49              & 0,55               & 0,40              & 373 \\
		\hline
		\textit{Acurácia}  & ---               & ---                & 0,40              & 373 \\
		\hline
	\end{tabular}
	\legend{Fonte: Elaborada pelo autor.}
\end{table}

\begin{table}[!htbp]
	\caption[Matriz de confusão — FinBERT-PT-BR]{Matriz de confusão do FinBERT-PT-BR. Linhas: rótulo verdadeiro; colunas: rótulo predito.}
	\label{tab:cm_finbert}
	\centering
	\begin{tabular}{lccc}
		\hline
		& \textbf{Pred. Positivo} & \textbf{Pred. Neutro} & \textbf{Pred. Negativo} \\
		\hline
		\textbf{Real Positivo} & 28                      & 7                     & 24                      \\
		\hline
		\textbf{Real Neutro}   & 39                      & 81                    & 146                     \\
		\hline
		\textbf{Real Negativo} & 2                       & 4                     & 42                      \\
		\hline
	\end{tabular}
	\legend{Fonte: Elaborada pelo autor.}
\end{table}

O FinBERT-PT-BR alcançou acurácia de 40,48\% e F1-score macro de 40,44\%, um resultado superior às abordagens léxicas, mas distante do desempenho documentado pelos autores do modelo sobre o corpus de notícias financeiras de treinamento \cite{santos2023finbert}. O padrão mais relevante na \autoref{tab:cm_finbert} é o viés pronunciado em favor da classe negativo: o modelo classificou 146 das 266 publicações neutras (54,9\%) e 24 das 59 positivas (40,7\%) como negativo, atingindo revocação de 88\% nessa classe à custa de precisão de apenas 20\%. Esse comportamento indica que o modelo generaliza excessivamente o padrão negativo aprendido a partir do corpus de notícias formais --- no qual termos como ``queda'', ``perda'' e ``risco'' são indicadores frequentes de sentimento negativo --- para publicações de formato editorial curto do X/Twitter, nas quais o mesmo vocabulário aparece recorrentemente em contextos factuais e neutros.

% ── 4.2.4 BERTimbau ──────────────────────────────────────────────────────────

\subsection{BERTimbau Ajustado}
\label{subsec:res_bertimbau}

O BERTimbau (\texttt{neuralmind/bert-base-portuguese-cased}) \cite{souza2020} foi submetido a \textit{fine-tuning} sobre a partição de treinamento do corpus rotulado, conforme o protocolo descrito na \autoref{subsubsec:protocolo_inferencia_bertimbau}. Os resultados, apresentados na \autoref{tab:metrics_bertimbau}, revelam desempenho amplamente superior às demais abordagens avaliadas.

\begin{table}[!htbp]
	\caption[Métricas de desempenho — BERTimbau ajustado]{Métricas de desempenho do BERTimbau ajustado ($n = 373$)}
	\label{tab:metrics_bertimbau}
	\centering
	\begin{tabular}{lcccc}
		\hline
		\textbf{Classe}    & \textbf{Precisão} & \textbf{Revocação} & \textbf{F1-Score} & \textbf{Suporte} \\
		\hline
		Positivo           & 0,84              & 0,80               & 0,82              & 59  \\
		\hline
		Neutro             & 0,91              & 0,97               & 0,94              & 266 \\
		\hline
		Negativo           & 0,88              & 0,62               & 0,73              & 48  \\
		\hline
		\textit{Macro avg} & 0,88              & 0,80               & 0,83              & 373 \\
		\hline
		\textit{Acurácia}  & ---               & ---                & 0,90              & 373 \\
		\hline
	\end{tabular}
	\legend{Fonte: Elaborada pelo autor.}
\end{table}


\begin{table}[!htbp]
	\caption[Matriz de confusão — BERTimbau ajustado]{Matriz de confusão do BERTimbau ajustado. Linhas: rótulo verdadeiro; colunas: rótulo predito.}
	\label{tab:cm_bertimbau}
	\centering
	\begin{tabular}{lccc}
		\hline
		& \textbf{Pred. Positivo} & \textbf{Pred. Neutro} & \textbf{Pred. Negativo} \\
		\hline
		\textbf{Real Positivo} & 47                      & 10                    & 2                       \\
		\hline
		\textbf{Real Neutro}   & 7                       & 257                   & 2                       \\
		\hline
		\textbf{Real Negativo} & 2                       & 16                    & 30                      \\
		\hline
	\end{tabular}
	\legend{Fonte: Elaborada pelo autor.}
\end{table}

O BERTimbau ajustado alcançou acurácia de 89,54\% e F1-score macro de 82,84\%, representando ganhos absolutos de 49,1 pontos percentuais em acurácia e 42,4 pontos em F1-score macro em relação ao FinBERT-PT-BR. A matriz de confusão (\autoref{tab:cm_bertimbau}) evidencia que o modelo distribuiu os erros de forma muito mais equilibrada entre as três classes: 257 das 266 instâncias neutras (96,6\%) foram corretamente classificadas, 47 das 59 positivas (79,7\%) e 30 das 48 negativas (62,5\%). Diferentemente do FinBERT-PT-BR, o BERTimbau ajustado não apresenta viés sistemático em favor de nenhuma classe; os erros residuais concentram-se principalmente em confusão entre neutro e as classes minoritárias (10 positivos classificados como neutro e 16 negativos como neutro), um padrão esperado dado o desequilíbrio de classes no corpus.
% ============================================================
%  USPSC-Cap4-Avaliacao_Experimental_4.3-Analise_Comparativa.tex
%  Seção 4.3 — Análise Comparativa
% ============================================================

\section{Análise Comparativa}
\label{sec:analise_comparativa}

A \autoref{tab:comparativa_geral} consolida os resultados das quatro abordagens avaliadas sobre o mesmo conjunto de teste ($n = 373$), permitindo a comparação direta de acurácia, F1-score por classe e F1-score macro.

\begin{table}[!htbp]
	\caption[Comparativo de desempenho entre as abordagens]{Comparativo de desempenho entre as quatro abordagens ($n = 373$). F1-scores expressos por classe (Pos, Neu, Neg) e como média macro (Macro). Melhor resultado por coluna em negrito.}
	\label{tab:comparativa_geral}
	\centering
	\begin{tabular}{lccccc}
		\hline
		\textbf{Abordagem}  & \textbf{Acurácia} & \textbf{F1 Pos} & \textbf{F1 Neu} & \textbf{F1 Neg} & \textbf{F1 Macro} \\
		\hline
		OpLexicon v3.0      & 0,26              & 0,29            & 0,22            & 0,28            & 0,26              \\
		\hline
		SentiLex-PT         & 0,29              & 0,29            & 0,31            & 0,27            & 0,29              \\
		\hline
		FinBERT-PT-BR       & 0,40              & 0,44            & 0,45            & 0,32            & 0,40              \\
		\hline
		BERTimbau (FT)      & \textbf{0,90}     & \textbf{0,82}   & \textbf{0,94}   & \textbf{0,73}   & \textbf{0,83}     \\
		\hline
	\end{tabular}
	\legend{Fonte: Elaborada pelo autor. (FT) = \textit{fine-tuned} sobre o corpus local.}
\end{table}


\subsection{Hierarquia de Desempenho}
\label{subsec:hierarquia}

Os resultados estabelecem uma hierarquia clara entre as abordagens, ordenada por complexidade metodológica crescente: OpLexicon < SentiLex-PT < FinBERT-PT-BR $\ll$ BERTimbau ajustado. O salto de desempenho mais expressivo ocorre entre o FinBERT-PT-BR e o BERTimbau ajustado: a adaptação ao gênero e ao subdomínio via \textit{fine-tuning} local produziu ganhos de 49,1 pontos percentuais em acurácia e 42,4 pontos em F1-score macro, valores que colocam o BERTimbau ajustado como a única abordagem com desempenho praticamente utilizável para classificação de sentimento neste corpus.

Entre as abordagens léxicas, o SentiLex-PT supera marginalmente o OpLexicon em F1-score macro (29,00\% versus 26,12\%), diferença atribuível, em parte, à maior cobertura morfológica do SentiLex-PT, que inclui entradas anotadas com classes gramaticais e permite maior precisão na recuperação de termos flexionados. Nenhuma das duas abordagens, contudo, supera o limiar de desempenho de um classificador por maioria de classe --- que obteria acurácia de 71,3\% ao classificar todas as instâncias como neutro --- evidenciando que a tarefa não pode ser resolvida por correspondência léxica simples neste tipo de corpus.

\subsection{Desempenho nas Classes Minoritárias}
\label{subsec:classes_minoritarias}

O desempenho nas classes minoritárias (positivo e negativo) é o discriminador mais revelador entre as abordagens. A \autoref{tab:comparativa_minoritarias} detalha precisão e revocação das três classes para cada abordagem.

\begin{table}[!htbp]
	\caption[Precisão e revocação por classe]{Precisão (P) e revocação (R) por classe para cada abordagem}
	\label{tab:comparativa_minoritarias}
	\centering
	\begin{tabular}{lcccccc}
		\hline
		& \multicolumn{2}{c}{\textbf{Positivo}} & \multicolumn{2}{c}{\textbf{Neutro}} & \multicolumn{2}{c}{\textbf{Negativo}} \\
		\textbf{Abordagem} & \textbf{P}   & \textbf{R}              & \textbf{P}  & \textbf{R}            & \textbf{P}   & \textbf{R}              \\
		\hline
		OpLexicon v3.0     & 0,18         & 0,68                    & 0,81        & 0,13                  & 0,20         & 0,46                    \\
		\hline
		SentiLex-PT        & 0,19         & 0,59                    & 0,78        & 0,20                  & 0,18         & 0,48                    \\
		\hline
		FinBERT-PT-BR      & 0,41         & 0,47                    & 0,88        & 0,30                  & 0,20         & 0,88                    \\
		\hline
		BERTimbau (FT)     & 0,84         & 0,80                    & 0,91        & 0,97                  & 0,88         & 0,62                    \\
		\hline
	\end{tabular}
	\legend{Fonte: Elaborada pelo autor.}
\end{table}

As abordagens léxicas apresentam revocação aparentemente elevada para a classe positivo (OpLexicon: 68\%; SentiLex-PT: 59\%), mas com precisão inferior a 20\%, o que indica que a maioria das predições positivas é incorreta --- a proporção de verdadeiros positivos no total de predições positivas é baixíssima. Esse padrão reflete a tendência das abordagens léxicas de superclassificar instâncias como positivo quando o vocabulário não apresenta polaridade negativa explícita, mesmo que o conteúdo seja factual e neutro.

O FinBERT-PT-BR apresenta o inverso para a classe negativo: revocação de 88\% com precisão de apenas 20\%, confirmando o viés pronunciado discutido na \autoref{subsec:res_finbert}. O modelo recupera quase todas as instâncias verdadeiramente negativas, mas ao custo de classificar incorretamente grande parte das neutras como negativas.

O BERTimbau ajustado é a única abordagem que simultaneamente mantém precisão e revocação acima de 60\% para todas as três classes, incluindo as minoritárias. Esse resultado indica que o \textit{fine-tuning} local capacitou o modelo a discriminar os padrões linguísticos específicos das três classes no contexto de publicações financeiras do X/Twitter em português, sem privilegiar sistematicamente nenhuma das classes ao custo das demais.


\subsection{Viés de Classe Neutra}
\label{subsec:vies_neutro}

O tratamento da classe neutra é um indicador relevante da qualidade de cada abordagem para a tarefa específica deste trabalho. Publicações jornalísticas financeiras são predominantemente neutras --- 71,3\% do corpus de teste --- e a correta identificação dessa classe é condição para resultados de classificação de sentimento fiéis à distribuição real de polaridade do corpus. A \autoref{tab:comparativa_geral} mostra que somente o BERTimbau ajustado alcança F1-score acima de 0,80 para a classe neutra (F1 = 0,94). As demais abordagens variam entre 0,22 (OpLexicon) e 0,45 (FinBERT-PT-BR) para essa classe, o que implica que essas abordagens subestimariam sistematicamente a proporção de publicações neutras e exagerariam a fração de publicações com polaridade definida.


% ============================================================
%  USPSC-Cap4-Avaliacao_Experimental_4.4-Discussao.tex
%  Seção 4.4 — Discussão
% ============================================================

\section{Discussão}
\label{sec:discussao_resultados}

\subsection{Especialização Prévia versus Ajuste Local ao Gênero}
\label{subsec:discussao_especializacao}

O contraste mais relevante que emerge dos resultados é o que opõe o FinBERT-PT-BR ao BERTimbau ajustado, duas abordagens neurais que diferem fundamentalmente na estratégia de especialização ao domínio e ao gênero. O FinBERT-PT-BR incorpora \textbf{especialização prévia de domínio} (\textit{prior domain specialization}): o modelo foi pré-treinado sobre um corpus de notícias financeiras formais em português \cite{santos2023finbert}, o que lhe confere vocabulário financeiro robusto, mas não o prepara para as particularidades do gênero textual de publicações do X/Twitter --- sentenças curtas, uso de emojis convertidos, menções a tickers, estilo telegráfico e ausência de marcadores explícitos de polaridade. O BERTimbau ajustado incorpora \textbf{especialização pós-hoc de gênero e subdomínio} (\textit{post hoc genre and subdomain specialization}): o modelo base generalista \cite{souza2020} foi adaptado diretamente sobre publicações do mesmo gênero e subdomínio que compõem o corpus de avaliação.

Os resultados confirmam a hipótese de que a adaptação local ao gênero produz ganhos substanciais mesmo quando comparada a um modelo com especialização de domínio prévia: o BERTimbau ajustado superou o FinBERT-PT-BR em 34,2 pontos percentuais de F1-score macro (77,45\% versus 43,20\%). Esse achado está alinhado com a tradição de \textit{fine-tuning} de modelos de linguagem em domínios especializados \cite{devlin2019,santos2023finbert}, que demonstra que a proximidade entre os dados de ajuste e os dados de inferência é um fator determinante para o desempenho em tarefas de classificação de sentimento.

\subsection{Diagnóstico do Viés do FinBERT-PT-BR}
\label{subsec:discussao_vies_finbert}

O viés do FinBERT-PT-BR em favor da classe negativo (revocação de 88\%, precisão de 20\%) merece análise específica. Uma hipótese explicativa plausível é a distribuição assimétrica de sentimento no corpus de treinamento do modelo: notícias financeiras formais tendem a tratar de eventos de risco, quedas de indicadores, inadimplência e crises, gerando uma representação mais densa de vocabulário negativo no espaço de \textit{embeddings} do modelo. Quando aplicado a publicações do X/Twitter, o modelo interpreta o vocabulário neutro de cobertura de mercado (e.g., ``variação'', ``fechamento'', ``resultado'') como sinal negativo, produzindo o padrão de superclassificação observado. Esse diagnóstico é consistente com a observação de \citeonline{loughran2011}, que documentaram que léxicos de sentimento geral aplicados ao domínio financeiro frequentemente atribuem polaridade incorreta a termos com significado neutro ou positivo em contextos de negócios.

\subsection{Limitações das Abordagens Léxicas em Corpus de Tweet}
\label{subsec:discussao_lexical}

O desempenho das abordagens léxicas --- OpLexicon (F1 macro 26,83\%) e SentiLex-PT (F1 macro 30,79\%) --- é sistematicamente inferior ao que seria esperado de uma aplicação ao domínio para o qual foram construídos: textos em português com marcadores explícitos de polaridade. Dois fatores estruturais explicam esse resultado. Primeiro, publicações jornalísticas financeiras do X/Twitter são predominantemente informativas e factuais; a proporção de termos com polaridade lexical explícita é baixa, o que reduz a capacidade de discriminação dos léxicos. Segundo, a ausência de mecanismos de desambiguação contextual implica que termos polissêmicos ou dependentes de contexto (e.g., ``alta'' como variação de preço versus expressão de intensidade) são classificados com base apenas em sua ocorrência, independentemente do sentido contextual \cite{liu2020survey}.

Ressalta-se que o OpLexicon e o SentiLex-PT foram construídos para o português de domínio geral, sem especialização para o vocabulário financeiro; o desempenho observado pode ser interpretado como um \textit{baseline} léxico de referência, não como o limite superior das abordagens baseadas em léxico para o domínio.


\subsection{Restrições à Generalização dos Resultados}
\label{subsec:discussao_restricoes}

Os resultados devem ser interpretados considerando as restrições metodológicas identificadas na \autoref{sec:limitacoes}. O conjunto de teste foi anotado por um único avaliador, inviabilizando o cálculo de concordância interanotadores; a margem de incerteza associada à qualidade do \textit{gold standard} não é quantificada. Adicionalmente, o corpus é restrito a dois veículos de mídia financeira institucional (InfoMoney e InvestingBrasil), o que circunscreve a representatividade dos resultados ao gênero editorial financeiro em língua portuguesa --- os achados podem não ser diretamente generalizáveis a outros gêneros textuais do domínio financeiro, como análises de analistas independentes, fóruns de investidores ou mensagens de corretoras.

Por fim, o tamanho do corpus de ajuste do BERTimbau (aproximadamente 2.600 instâncias de treinamento) é modesto em comparação com os corpora utilizados em estudos de \textit{fine-tuning} de grande escala; a utilização de conjuntos de dados maiores e de maior diversidade de gênero poderia produzir modelos com desempenho superior nas classes minoritárias, especialmente positivo (F1 = 0,72) e negativo (F1 = 0,69).